%----------------------------------------------------------------------------
\chapter{Archiving the code on Zenodo}\label{app:zenodo}
%----------------------------------------------------------------------------
A reviewer assessing \textbf{correctness} (criterion D/2) needs the code behind your results. A GitHub link can change or disappear.
\href{https://zenodo.org/}{Zenodo}, run by CERN, permanently archives GitHub releases and assigns DOIs. Cite the archive.
These steps follow Zenodo's September 2026 documentation. If its interface has changed, follow the linked pages.

\section{Prepare the repository}
Prepare the \texttt{code/} submodule for archiving (\cref{app:code:github}); Zenodo archives this repository, not the notes repository.
\begin{enumerate}
  \item Complete \texttt{README.md} using \cref{app:code:readme}.
  \item Add a \texttt{LICENSE} file and ask us which license to use: Apache 2.0 for tools (\cref{app:code:readme}), MIT otherwise.
  \item Add \texttt{.zenodo.json} at the repository root (\cref{lst:zenodo-json}). It supplies the record's title, authors, license and links, and lists your supervisors as \texttt{contributors} of type \texttt{Supervisor}; see the \href{https://help.zenodo.org/docs/github/describe-software/zenodo-json/}{format and example}. The code template's \texttt{hygiene} workflow checks that it parses and that \texttt{CITATION.cff} validates.
  \item Link related works with \texttt{related\_identifiers} (\cref{lst:zenodo-related}): \texttt{isDerivedFrom} the upstream repository of a fork, and \texttt{isSupplementTo} the paper once it has a DOI. Leave the field out until you have a real identifier, because a placeholder fails the release.
  \item If the repository has a \texttt{CITATION.cff}, keep its fields in line with \texttt{.zenodo.json}; Zenodo uses \texttt{.zenodo.json} when both exist.
\end{enumerate}

\begin{listing}[!ht]
\begin{minted}{json}
{
  "title": "TODO Artifact for: <thesis title>",
  "description": "Code and measurement scripts for the TDK paper <thesis title>.",
  "upload_type": "software",
  "version": "1.0.0",
  "publication_date": "TODO YYYY-MM-DD",
  "license": "MIT",
  "creators": [
    { "name": "TODO Surname, Given name", "affiliation": "Budapest University of Technology and Economics", "orcid": "TODO your ORCID iD, or remove this field" }
  ],
  "contributors": [
    { "name": "Farkas, Martin", "type": "Supervisor", "affiliation": "Budapest University of Technology and Economics" },
    { "name": "Kocsis, Imre", "type": "Supervisor", "affiliation": "Budapest University of Technology and Economics" }
  ],
  "keywords": ["TODO"]
}
\end{minted}
\caption{A minimal \texttt{.zenodo.json}.}\label{lst:zenodo-json}
\end{listing}

\begin{listing}[!ht]
\begin{minted}{json}
"related_identifiers": [
  { "identifier": "https://github.com/<upstream>/<repo>", "relation": "isDerivedFrom", "resource_type": "software" },
  { "identifier": "10.xxxx/<paper DOI>", "relation": "isSupplementTo", "resource_type": "publication-article" }
]
\end{minted}
\caption{Related works in \texttt{.zenodo.json}: the upstream of a fork and the paper.}\label{lst:zenodo-related}
\end{listing}

\section{Before the repository goes public}\label{app:zenodo:public}
Making the repository public, on GitHub or through its Zenodo archive, publishes its whole history, not only the current files. A file you deleted last month is still in a commit, and Zenodo does not let you change the files of a published record. We do not rewrite the history: keep the current files clean, and review what the history still holds. Check the following before the first release:
\begin{itemize}
  \item \textbf{No links to private or unpublished places.} No Overleaf project, paper or notes repository, Notion, Teams or SharePoint page, or chat share link, in the README or anywhere else. Link only public URLs; the same-year rule for companion papers (\cref{app:tdk:companion}) applies here too.
  \item \textbf{No paper sources.} The paper lives in the paper repository; standalone TikZ figures that a script generates may stay.
  \item \textbf{You author your commits.} Configure the agent to commit as you; an agent identity such as \texttt{Claude <noreply@anthropic.com>} lists the agent as a contributor on GitHub. A \texttt{Co-Authored-By} line naming the agent is fine, and so are commits by your advisors. Do not add an AI app or bot as a collaborator.
  \item \textbf{No agent files in the archive.} The code template's \texttt{.gitattributes} marks \path{.claude/}, \texttt{CLAUDE.md}, \texttt{AGENTS.md} and similar files \texttt{export-ignore}, which keeps them out of GitHub's release archive that Zenodo stores. Add any other agent file there. If your \texttt{CLAUDE.md} is really documentation of the code, rename it to \texttt{README.md} or move it under \path{docs/}; otherwise the archive silently drops it.
  \item \textbf{No secrets, test keys included.} Keep keys, wallets and tokens in an untracked \texttt{.env}, and commit an \texttt{.env.example} with placeholders. A key that was ever committed is burned, because the history keeps it: rotate it, so that the committed value no longer works or holds anything.
  \item \textbf{No machine-specific paths} such as \path{/home/you/}: use relative paths or variables.
  \item \textbf{A fork credits its upstream.} Create it with GitHub's \emph{Fork} button from the start: GitHub cannot turn an existing repository into a fork later, and only a real fork shows the ``forked from'' link. Keep the upstream license, name its authors in \texttt{.zenodo.json}, link the upstream with \texttt{isDerivedFrom} (\cref{lst:zenodo-related}), and say in the README what you changed.
\end{itemize}
The code template's \texttt{hygiene} workflow checks these on every push: secrets with gitleaks over the whole history, the rest with \texttt{scripts/hygiene.py}. Findings in the current files fail the run. Findings in the history only warn, because the history stays as it is: read each one, and decide whether it is harmless once public. A secret in the history fails the run until you rotate it; then add the finding's exact fingerprint to \texttt{.gitleaksignore}, with a one-line reason such as \enquote{rotated on 2026-11-02} directly above it. Use the same file for findings that are not secrets, such as a public key in an example. \texttt{just hygiene --strict}, which fails on history findings too, remains an optional local check.

\paragraph{The publication gate.} When you start preparing the release, copy \path{scripts/publishing-checklist.md} to \texttt{PUBLISHING.md}. From then on, the workflow runs \texttt{just hygiene --publish} instead: it also fails on unfinished Zenodo metadata, missing supervisors, a README with \texttt{TODO}s or without its sections (\cref{app:code:readme}), missing lockfiles, and a missing release workflow. \texttt{PUBLISHING.md} lists what a script cannot check, such as rotated secrets, reviewed history warnings and the license choice. Fill in the evidence for each item; we sign it off by approving your pull request. Release only when the gate passes.

\section{Publish the first release}
\begin{enumerate}
  \item Sign in to \href{https://zenodo.org/}{Zenodo} with the repository owner's GitHub account. In the profile menu, choose \emph{GitHub}, click \emph{Sync now}, find the repository and enable its toggle; see \href{https://help.zenodo.org/docs/github/enable-repository/}{enable a repository}.
  \item Commit the prepared files. In GitHub's \emph{Releases}, choose \emph{Draft a new release}. Create a tag such as \texttt{v1.0.0} at the commit used for the measurements. Add release notes and publish.
  \item The code template's \texttt{release} workflow then pushes the runtime image as \texttt{ghcr.io/<you>/<repo>:v1.0.0}. Copy \path{.github/workflows/release.yml} if your repository did not start from the template. A new package is \href{https://docs.github.com/en/packages/learn-github-packages/configuring-a-packages-access-control-and-visibility}{private by default}, even when the repository is public: in the package's settings, change its visibility to public, or the README's \texttt{docker run} fails for reviewers.
  \item Open the repository on Zenodo's GitHub page and wait for processing. Check the new record's files and metadata. A failed release has an \emph{Errors} section; see \href{https://help.zenodo.org/docs/github/archive-software/github-upload/}{archive a release}.
\end{enumerate}
If you cannot connect the repository, upload an archive manually. Choose \emph{+}, then \emph{New upload}, fill in metadata, select \emph{Preview}, then \emph{Publish}; see the \href{https://help.zenodo.org/docs/deposit/create-new-upload/}{upload instructions}.

\section{Cite it}\label{app:zenodo:cite}
Each release has a \emph{version DOI}. The \emph{concept DOI} identifies all versions and resolves to the newest; see \href{https://support.zenodo.org/help/en-gb/1-upload-deposit/97-what-is-doi-versioning}{DOI versioning}.
\begin{itemize}
  \item Cite the \textbf{concept DOI} in the paper, following the \texttt{artifact} entry in \texttt{bib/references.bib}, in \tplref{ch:implementation}. Beside it, name the \textbf{results tag}: the Git tag and corresponding Zenodo version used for every reported number. The DOI leads to the latest packaging; the tag identifies the exact measured version.
  \item Put the concept DOI badge in \texttt{README.md} (\cref{app:code:readme}). Copy the badge's Markdown from the record page's DOI badge dialog.
  \item Copy both DOIs from the Zenodo record; never guess a DOI by changing its number.
\end{itemize}

\section{Release a new version}\label{app:zenodo:versions}
After fixing code, update \texttt{version} and \texttt{publication\_date} in \texttt{.zenodo.json}, commit, and publish a release under a new tag, such as \texttt{v1.0.1}. Zenodo adds a version with its own DOI to the record; see \href{https://help.zenodo.org/docs/deposit/manage-versions/}{manage versions}. Never move an old tag or rerun the paper's measurements against a version it does not describe. Keep the paper's results tag unchanged.

The lifecycle (\cref{tab:lifecycle}) allows code polish \emph{after} the paper deadline, because the paper cites the concept DOI, which opens the version current when a reviewer visits. You have about two days after upload to finish the README, the Docker images, the results map and the reproduction-limits section, and to release them as a new version. Polish after that is allowed, but reviewers may already have looked.
